% Generated draft: requires mathematical and source review.
\documentclass[11pt]{article}
\usepackage[margin=25mm]{geometry}
\usepackage{fontspec}
\setmainfont{Latin Modern Roman}
\usepackage{amsmath,amssymb,mathtools,unicode-math}
\setmathfont{Latin Modern Math}
\usepackage{enumitem,needspace,xurl,hyperref}
\setlength{\parindent}{0pt}
\setlength{\parskip}{5pt}
\setlength{\emergencystretch}{3em}
\newcommand{\N}{\mathbb{N}}
\newcommand{\Z}{\mathbb{Z}}
\newcommand{\Q}{\mathbb{Q}}
\newcommand{\R}{\mathbb{R}}
\newcommand{\C}{\mathbb{C}}
\newcommand{\F}{\mathbb{F}}
\newcommand{\A}{\mathbb{A}}
\newcommand{\PP}{\mathbb{P}}
\newcommand{\E}{\mathbb{E}}
\newcommand{\eps}{\varepsilon}
\newcommand{\dd}{\,\mathrm{d}}
\begin{document}
\section*{UnsolvedMath ID: 30000666}
\textbf{Source problem code:} catalogue entry
\subsection{Definitions and mathematical statement}
Let \(k\) be the source's base field, let \(G\) be a connected semisimple algebraic group over \(k\), and let \(\rho:G\to\operatorname{GL}(V)\) be a finite-dimensional algebraic representation with trivial stabilizers in general position. In the standard geometric interpretation this means that stabilizers are trivial on a dense open subset after the intended base change. The field, characteristic, and scheme-theoretic versus pointwise stabilizer convention are not supplied.

For a \(k\)-defined algebraic automorphism \(\sigma:G\to G\), use the twist convention
\[
g\star v=\rho(\sigma(g))v\qquad\text{on }V^{(\sigma)}.
\]
The general question is whether there exists a birational map with rational inverse
\[
\Phi:V\dashrightarrow V^{(\sigma)},\qquad
\Phi(\rho(g)v)=\rho(\sigma(g))\Phi(v)
\]
where the expressions are defined, for every indicated \(G,V,\sigma\). This is birational \(G\)-isomorphism of the actual representations, not merely stable equivalence after adding trivial factors or an isomorphism of their invariant fields. Some sources precompose by \(\sigma^{-1}\); the chosen convention must be matched to the report.

The separate explicit question, identified as Problem 4 by supplied verification, takes
\[
G=\operatorname{SL}_d/\mu_d,\qquad V=\operatorname{Sym}^d(k^d),\qquad
\sigma([g])=[(g^{\mathsf T})^{-1}].
\]
Here \(\mu_d\) is the central subgroup of scalar \(d\)-th roots of unity. Its action on \(\operatorname{Sym}^d(k^d)\) is trivial, giving the stated quotient-group action. Is this \(V\) birationally \(G\)-isomorphic to its specified inverse-transpose twist? The record reports trivial stabilizers in general position for \(d>3\) in the intended field setting, making that range an instance of the general question. It does not thereby assert the same stabilizer property for \(d\leq3\), nor exclude those dimensions from the separately posed specialization without a source range.

The supplied literature triage checked 2026-08-21 reports no authoritative complete resolution of the extracted pair of problems. No additional field hypotheses or partial theorem are given, so the generic-stabilizer assertion is retained as supplied rather than certified uniformly in every characteristic.
\subsection{Short English statement}
Ask whether a generically free representation of a connected semisimple group is birationally equivariantly isomorphic to every automorphism twist, especially the stated symmetric-power example with inverse-transpose twist. The base-field and stabilizer conventions remain essential.
\subsection{Sources}
\begin{itemize}
\item Local UnsolvedMath registry, permanent ID 30000666, catalogue entry, \emph{Birational Equivalence of Twisted Algebraic-Group Representations}.
\item \emph{Affine Algebraic Geometry} (2007), Oberwolfach Reports, pp. 5--82, DOI \url{https://doi.org/10.4171/owr/2007/01}; supplied verification identifies the general question as Problem 3 and specialization as Problem 4.
\item Supplied literature triage checked 2026-08-21 cites only that report; no separate theorem or precise field assumptions for generic freeness are supplied.
\end{itemize}
\end{document}
